\section{Institutional background}\label{sec:institutional}

\subsection{The Peruvian minimum wage}

Peru sets a single national minimum wage, the \emph{Remuneración Mínima Vital}
(RMV), which applies to workers in the private-sector labor regime. Unlike federal
systems such as the United States or Brazil, and unlike countries with occupational
or regional wage floors, Peru has no sub-national or industry-specific statutory
minimum in the general regime. A small number of activities carry fixed multiples of
the RMV (for example the mining premium and the agricultural daily wage), but these
are pegged mechanically to the national value and move with it.\footnote{These pegged
regimes, such as the agrarian and mining regimes and the small-enterprise (MYPE) regime,
move one-for-one with the RMV and so introduce no cross-sectional variation in its level.} The consequence for
research design is decisive: because the level of the RMV is common to the whole
country, there is no cross-sectional variation in the policy and no staggered timing
of adoption. Any credible research design must therefore exploit differences in how
exposed workers are to a change that reaches everyone simultaneously.

In principle the RMV is to be adjusted periodically by the tripartite National Labor
Council in line with inflation and productivity. In practice the Council has never
set the RMV, and the level is fixed at the discretion of the executive through a
Supreme Decree. Adjustments are consequently infrequent and irregular in both timing
and magnitude, and Peru is among the most volatile minimum wage setters in the
region \citep{jimenez_2023, cespedes_2005}. The RMV had stood at 930 soles since
April 2018.\footnote{Fixed by Supreme Decree 004-2018-TR. The two preceding increases raised
the RMV from 750 to 850 soles in May 2016 and from 850 to 930 soles in April 2018, so the
2022 reform followed an unusually long freeze.}

Two features make the RMV unusually binding by international standards. First, its
ratio to the median wage is high. In my pre-reform sample the national median
monthly wage of a dependent employee is almost exactly 930 soles, so the old minimum
already sat at the median of the covered wage distribution and the Kaitz index, the
ratio of the minimum to the median wage, is close to one nationally and well above
one in the poorer departments. Second, compliance is far from universal. Consistent
with earlier work \citep{jaramillo_lopez_2006, cespedes_2005}, a large share of
workers earn below the statutory floor, especially in rural areas, in small firms,
and in the informal sector, while compliance inside registered formal firms is high.
Enforcement capacity is limited.\footnote{Labor inspection is the responsibility of the
National Superintendency of Labor Inspection (SUNAFIL), whose inspection staff is thin
outside Lima. Earlier work estimates overall non-compliance with the RMV at roughly 30 to
40 percent, concentrated in rural areas, small firms, and the informal sector
\citep{jaramillo_lopez_2006, cespedes_2005}.} These features imply that a nominal increase in the
RMV can bind mechanically for formal workers near the floor while leaving much of the
informal sector formally uncovered, so that the policy's real reach depends on the
degree of informality and on the ease of movement across the formal and informal
margins.

\subsection{The 2022 reform}

The reform I study was enacted by Supreme Decree 003-2022-TR, signed on 2 April 2022
and published on 3 April, and took effect on 1 May 2022. It raised the RMV from 930
to 1,025 soles, an increase of 95 soles or 10.2 percent, the first change in four
years. Figure~\ref{fig:mw} plots the nominal and real value of the RMV over
2021--2025. Two points are worth noting. The nominal increase of 10.2 percent was
partly offset by inflation, which ran at roughly 8 percent over 2022, so the increase
in the real value of the floor was more modest and eroded over the following two
years. This tempers the expected size of any wage response and is relevant when
comparing my estimates with those from lower-inflation episodes. The reform also
occurred during the post-pandemic recovery, a period of unusual labor market
turbulence that I address explicitly in the empirical analysis.

A second reform, enacted by Supreme Decree 006-2024-TR and effective on 1 January
2025, raised the RMV further to 1,130 soles. This change falls just outside my main
study window, but the ENAHO quarters for 2025 are available and I use this second
reform as an out-of-sample replication of the design in Section~\ref{sec:robustness}.
The two reforms are of very similar magnitude, 10.2 percent each, which makes the
2025 episode a natural validation exercise in a calmer, post-pandemic environment.

\begin{figure}[H]
\centering
\includegraphics[width=0.82\textwidth]{fig1_minimum_wage.pdf}
\caption{The nominal and real Peruvian minimum wage, 2021--2025. The real series is
deflated by the Lima consumer price index (December 2021 = 100). Vertical dotted
lines mark the reforms effective 1 May 2022 (Supreme Decree 003-2022-TR) and 1
January 2025 (Supreme Decree 006-2024-TR).}
\label{fig:mw}
\end{figure}
